\documentclass[12pt,a4paper,titlepage]{article}
\usepackage[latin1]{inputenc}
\usepackage[brazil]{babel}

\usepackage[T1]{fontenc} % permite hifeniza��o em palavras acentuadas
\usepackage[scaled]{helvet} % usa uma fonte sem serifa no tamanho adequado
\renewcommand*\familydefault{\sfdefault} % usa a fonte n�o serifada

\usepackage{ae}
\usepackage{graphicx} % para poder inserir imagens PNG, JPG e GIG
\usepackage{verbatim}
\usepackage{rotating}
\usepackage{fullpage}
\usepackage{indentfirst}
\usepackage{moreverb}
\newcommand{\LinhaHorizontal}{\rule{\linewidth}{0.5mm}}
\usepackage[show]{ed} % para adicionar notas de edi��o - retirar da vers�o final (usar o comando \ednote{})
\def\verbatimtabsize{2\relax}
\usepackage{color}
	\definecolor{contentsLinks}{RGB}{0, 0, 100}
\usepackage[breaklinks]{hyperref}
\hypersetup{
    pdfnewwindow=true,			% links in new window
    colorlinks=true,			% false: boxed links; true: colored links
    linkcolor=contentsLinks,	% color of internal links
    citecolor=green,			% color of links to bibliography
    filecolor=magenta,			% color of file links
    urlcolor=cyan				% color of external links
}


\usepackage{verbatim}
% \def\verbatimtabsize{2\relax}

\usepackage{multirow}

%C�DIGOS FONTE
\usepackage{listings} % para importa��o de c�digos fonte
\renewcommand{\lstlistingname}{C�digo} % defini��o visual dos c�digos fonte
\lstset{
	extendedchars=true,
	inputencoding=latin1,
	belowcaptionskip=5pt,
	numbers=left,
	stepnumber=1,
	firstnumber=1,
	numberstyle=\tiny,
	breaklines=true,
	frame=tb,
	basicstyle=\ttfamily\small,
	showstringspaces=false,
	mathescape,
	tabsize=3,
	% prebreak = \raisebox{0ex}[0ex][0ex]{\ensuremath{\hookleftarrow}},
	keywordstyle=\ttfamily\color[rgb]{0,0,1},
	commentstyle=\ttfamily\color[rgb]{0.133,0.545,0.133},
	stringstyle=\ttfamily\color[rgb]{0.627,0.126,0.941}
}

\newcommand{\inputpng}[4]{%
\begin{figure}[!htb]
  \centering
  \includegraphics[width=#4\textwidth]{#1.png}
  \caption{\it #3.}
  \label{fig:#2}
\end{figure}
}
\begin{document}
	% T�tulo
	\begin{titlepage}
\begin{center}

% Cabe�alho
\includegraphics{figuras/Header}
\vfill\

% T�tulo
\LinhaHorizontal\
\\[0.8cm]
\begin{Huge}\bfseries{Relat�rio Final}\end{Huge}\\[.2cm]
\begin{Large}PCS-2056 : Linguagens e Compiladores\end{Large}
\\[0.4cm]
\LinhaHorizontal\
\vfill\

% Autores e Professor
\begin{minipage}{.4\textwidth}
	\begin{flushleft}
		\emph{Autores:}\\
		\begin{large}Eduardo Russo\\Helio Kazuo Nagamachi\end{large}
	\end{flushleft}
\end{minipage}
\begin{minipage}{.4\textwidth}
	\begin{flushright}
		\emph{Professor:}\\
		\begin{large}Ricardo Rocha\end{large}
	\end{flushright}
\end{minipage}
\\[1cm]

% Data
\begin{large}\today\end{large}

\end{center}
\end{titlepage}
	% Sum�rio
	\tableofcontents
	\newpage
	
	%
	% Cap�tulos
	%

	
	\section{Objetivo}
		Este documento tem como objetivo descrever a constru��o do compilador de Lua -- uma linguagem de programa��o de extens�o -- para um ambiente de execu��o escrito em Java para a segunda prova da disciplina de \textbf{Linguagens e Compiladores}.
	\section{Gram�tica}
		A gram�tica da linguagem Lua apresentada para a prova em Wirth -- c�digo \ref{code:codigos_Gramatica} -- foi reduzida de forma a necessitar de apenas tr�s Aut�matos de Pilha Estruturado (APEs).
\lstinputlisting[tabsize = 4, caption={Gram�tica.}, label = {code:codigos_Gramatica}]{codigos/Gramatica}

Para a redu��o, foi utilizado um aplicativo desenvolvido em Python para facilitar o trabalho. Esta redu��o � apresentada no c�digo \ref{code:codigos_Gramatica-reduzida}

\lstinputlisting[tabsize = 4, caption={Gram�tica reduzida.}, label = {code:codigos_Gramatica-reduzida}]{codigos/Gramatica-reduzida}
	
	\section{An�lise l�xica}
		Para a an�lise l�xica, foram identificadas as palavras reservadas na linguagem apresentadas a seguir:

\begin{table}[htb!]
	\label{table:label}
	\centering
	\begin{tabular}{p{2cm} p{2cm} p{2cm} p{2cm} p{2cm} p{2cm}}
		\#			& - 	& + 		& * 	& / 	& $\wedge$	\\
					&		&			&		&		&			\\
		< 			& <= 	& >			& >= 	& == 	& $\sim$=	\\
					&		&			&		&		&			\\
		: 			& [ 	& ] 		& = 	& \{ 	& \}		\\
					&		&			&		&		&			\\
		false 		& true 	& return	& not 	& break & local		\\
					&		&			&		&		&			\\
		if 			& then 	& elseif 	& else	& while & repeat	\\
					&		&			&		&		&			\\
		\% 			& . 	& .. 		& ... 	& and 	& or		\\
					&		&			&		&		&			\\
		\& 			& ; 	& ( 		& ) 	& end 	& nil 		\\
					&		&			&		&		&			\\
		function	& for 	& in		& do 	&until				\\
	\end{tabular}
\end{table}

Considerando essas palavras reservadas, mas as poss�veis sequ�ncias de palavras, foram criados aut�matos para reconhecer e retornar os \emph{tokens}. Seus tipos s�o apresentados no c�digo \ref{code:codigos_TokenType}.

\lstinputlisting[language=Java, tabsize = 4, caption={\emph{Tipos de tokens definidos: palavra reservada, identificador, n�mero, final de arquivo e cadeia de palavras.}}, label = {code:codigos_TokenType}]{codigos/TokenType.java}


A constru��o do l�xico utilizou os c�digos apresentados na figura \ref{fig:figuras_lexico} e sua classe principal -- que chama as fun��es dos outros c�digos -- encontra-se no c�digo \ref{code:anallex}.

\inputpng{figuras/lexico}{figuras_lexico}{C�digos utilizados para construir o analisador l�xico}{1.0}

\lstinputlisting[language=Java, tabsize = 4, caption={\emph{Classe ``Analyzer.java'' do pacote ``lex'', controla o processo de an�lise l�xica.}}, label = {code:anallex}]{codigos/Analyzerlex.java}
	\section{Analise sint�tica}
		A an�lise Sint�tica se utiliza de um automato estruturado de pilha que � descrito por \cite{JJ}, estrutura que utiliza uma pilha e alguns automatos finitos para realizar a tarefa de an�lise sint�tica.

\subsection{Estrutura}
O analisador sint�tico possu� basicamente um automato estruturado de pilha e acesso � c�digos que permitem o parse de arquivos de configura��o dos automatos finitos.

O Automato estruturado de pilha, tem sua estrutura definida por :
\subsubsection{Automato Finito}
O automato finito utilizado para a an�lise sint�tica possu� um vetor de booleanos que denota os seus estados, se uma posi��o do vetor possuir o booleano verdadeiro, quer dizer que aquele estado � final. O automato tamb�m possu� um vetor com as transi��es do automato.

\subsubsection{Transi��o} 
A transi��o pode ser de dois tipos, uma transi��o normal , ou a chamada de outro automato. Se a transi��o for normal, ela possu� o pr�ximo estado e o token que a ativa. Se a transi��o for uma chamada, a transi��o cont�m o n�mero do automato a ser chamado e o estado de retorno, que � um estado no automato atual que ser� o estado corrente ap�s a execu��o do segundo automato.

\subsubsection{Pilha}
A pilha � a estrutura utilizada pelo Automato estruturado de pilha para armazenar a informa��o de automato e estado quando a transi��o de um automato a ser executada � a chamada de outro automato. Quando esse tipo de transi��o ocorre o automato corrente passa a ser o automato denotado pela transi��o e a pilha guarda qual o automato que realizou a chamada, e qual o estado que esse deve voltar ap�s a execu��o do segunto automato.


\subsection{Parse}
A gram�tica em nota��o de Wirth foi reduzida de forma a ficarmos com 3 automatos, um de programa, um de bloco de c�digo e o terceiro de express�es. A gram�tica foi transformada em aut�matos finitos determin�sticos m�nimos com o uso da ferramenta (COLOCAR O RADIANTFIRE). Um exemplo da sa�da para um dos automatos encontra-se na se��o de anexo, c�digo  REFERENCIA!  A primeira linha do arquivo � sempre igual, indicando que o estado inicial do automato � o estado zero, a segunda linha mostra quais s�o os estados finais e as outras linhas indicam as transi��es. Se o elemento da transi��o representar um token � uma transi��o normal, caso seja o nome de um automato, indica a chamada para o outro automato.

\subsection{An�lise Sint�tica}



 
	\section{Defini��o do ambiente de execu��o}
		O ambiente de execu��o � composto por uma cole��o de bibliotecas escritas em baixo n�vel, assembly da MVN para prover algumas funcionalidades b�sicas. S�o elas:

\begin{itemize}
	\item fun��o l�gica OU
	\item fun��o l�gica E
	\item fun��o comparativa >
	\item fun��o comparativa <
	\item fun��o comparativa >=
	\item fun��o comparativa <=
	\item fun��o comparativa ==
	\item fun��o comparativa !=
	\item fun��o l�gica !
	\item fun��o print de strings
\end{itemize}

A ideia � manter esses arquivos com essas fun��es em arquivos separados, pelo uso da biblioteca do montador, ligador e relocador, um �nico arquivo do programa principal pode ser gerado.

\subsection{Registros de Ativa��o}
A implementa��o da MVN n�o tem uma pilha em seu ambiente de execu��o, isso implica na impossibilidade de realizar chamadas de fun��o recursivamente. Um modo de fazer isso � utilizando-se de registros de ativa��o para contornar o problema.

	\section{Estrutura de dados e algoritmos do compilador}
		% O objetivo deste projeto � modelar e construir um compilador que utilize uma linguagem de programa��o imperativa.
	\section{Sem�ntica din�mica}
		Para a constru��o do analisador sem�ntico, foram usados os c�digos apresentados na figura \ref{fig:analissem}.

\inputpng{figuras/semantica}{analissem}{C�digos utilizados para construir o analisador sem�ntico}{1.0}

Cada uma dessas classes apresentam-se nos c�digos \ref{code:codigos_Semantic}, \ref{code:codigos_TokensList} e \ref{code:codigos_TokenStack}.

\lstinputlisting[language=Java, tabsize = 4, caption={\emph{Classe ``Semantic.java'' do pacote ``semantic'', controla o processo de an�lise sem�ntica.}}, label = {code:codigos_Semantic}]{codigos/Semantic.java}


\lstinputlisting[language=Java, tabsize = 4, caption={\emph{Classe ``TokenList.java'' do pacote ``semantic''.}}, label = {code:codigos_TokensList}]{codigos/TokensList.java}


\lstinputlisting[language=Java, tabsize = 4, caption={\emph{Classe ``TokenStack.java'' do pacote ``semantic''.}}, label = {code:codigos_TokenStack}]{codigos/TokenStack.java}

	\section{Integra��o das rotinas sem�nticas}
	 	A implementa��o do compilador permite ao usu�rio as funcionalidades de mostrar na tela um inteiro de que esteja contido em uma vari�vel, express�es aritim�ticas, considerando prioridade e express�es booleanas sem considera��o a respeito das express�es boolenas.

Infelizmente as bibliotecas do montador, relocador e lingador apresentam algum erro em sua implementa��o e acabam inserindo no c�digo bin�rio final no meio do arquivo a seguinte linha: 0000 0002 , um desvio incondicional para o endere�o 2 de mem�ria, arruinando os resultados.

	
	% Refer�ncias
	% \nocite{*} % mostrar todos, n�o s� os citados
	% \bibliographystyle{apalike} % estilo da ABNT
	% \bibliography{bibliografia} % arquivo bibtex
	% \addcontentsline{toc}{section}{Refer�ncias}
\end{document}